\chapter{Error-Correcting Codes — Core Definitions}
%%%%%%%%%%%%%%%%%%%%%%%%%%%%%%%%%%%%%%%%%%%%%%%%%%%%%%%%%%%%%%%%%%%%%%%%%%%%%%%
\section{Overview}

This file documents the foundational definitions for $q$-ary block codes over a finite field
alphabet $\alpha$.  Throughout, $\alpha$ is a type carrying instances
\texttt{[Fintype α]}, \texttt{[Field α]}, \texttt{[DecidableEq α]}, and \texttt{[Nonempty α]}.

%%%%%%%%%%%%%%%%%%%%%%%%%%%%%%%%%%%%%%%%%%%%%%%%%%%%%%%%%%%%%%%%%%%%%%%%%%%%%%%
\section{Basic objects: codewords and codes}

\subsection{Codewords}

\begin{definition}[Codeword]
\lean{ErrorCorrectingCodes.Codeword}
\label{ErrorCorrectingCodes.Codeword}
\leanok
\statementsource{philips}{
  pdf-24e94ad30862-p023-b007,
  pdf-24e94ad30862-p025-b003
}
\statementsource{IEEE-TIT.23.157}{pdf-abed77da6b7a-p001-b004}
A codeword of length $n$ over alphabet $\alpha$ is a function $c : \mathrm{Fin}\,n \to \alpha$.
\end{definition}

\begin{definition}[Pointwise operations]
\lean{ErrorCorrectingCodes.Codeword.add}
\label{ErrorCorrectingCodes.Codeword.add}
\lean{ErrorCorrectingCodes.Codeword.sub}
\label{ErrorCorrectingCodes.Codeword.sub}
\lean{ErrorCorrectingCodes.Codeword.zero}
\label{ErrorCorrectingCodes.Codeword.zero}
\leanok
Pointwise addition, subtraction, and the all-zero codeword.
\end{definition}

\subsection{Codes and linear codes}

\begin{definition}[Code]
\lean{ErrorCorrectingCodes.Codeword.Code}
\label{ErrorCorrectingCodes.Codeword.Code}
\leanok
\statementsource{ada-coding-book}{pdf-7251cb8e4bc5-p015-b006}
\statementsource{philips}{pdf-24e94ad30862-p025-b003}
\statementsource{IEEE-TIT.23.157}{pdf-abed77da6b7a-p001-b004}
A (block) code of length $n$ over $\alpha$ is a finite set (\texttt{Finset}) of codewords.
\end{definition}

\begin{definition}[Linear code via generator matrix]
\lean{ErrorCorrectingCodes.Codeword.Linear_Code}
\label{ErrorCorrectingCodes.Codeword.Linear_Code}
\leanok
\statementsource{ada-coding-book}{
  pdf-7251cb8e4bc5-p045-b006,
  pdf-7251cb8e4bc5-p028-b003
}
A code $C$ is linear with generator matrix $G$ (shape $n\times m$) if every message
$c':\mathrm{Fin}\,m\to\alpha$ maps to a codeword $Gc'\in C$, and every codeword in $C$
arises this way.
\end{definition}

\begin{definition}[Linear code, existential form]
\lean{ErrorCorrectingCodes.Codeword.Linear_Code'}
\label{ErrorCorrectingCodes.Codeword.Linear_Code'}
\leanok
\statementsource{ada-coding-book}{
  pdf-7251cb8e4bc5-p045-b006,
  pdf-7251cb8e4bc5-p028-b003
}
Existential version: there exists some $m$ and generator matrix $G$ witnessing linearity.
\end{definition}

%%%%%%%%%%%%%%%%%%%%%%%%%%%%%%%%%%%%%%%%%%%%%%%%%%%%%%%%%%%%%%%%%%%%%%%%%%%%%%%
\section{Distance, weight, and rate}

\begin{definition}[$q$-ary entropy]
\lean{ErrorCorrectingCodes.Codeword.qaryEntropy}
\label{ErrorCorrectingCodes.Codeword.qaryEntropy}
\leanok
\statementsource{ada-coding-book}{pdf-7251cb8e4bc5-p069-b005}
\statementsource{IEEE-TIT.23.157}{pdf-abed77da6b7a-p001-b005}
\[
H_q(p) \;=\; p\log_q(q-1) - p\log_q p - (1-p)\log_q(1-p).
\]
\end{definition}

\begin{definition}[Hamming distance]
\lean{ErrorCorrectingCodes.Codeword.hamming_distance}
\label{ErrorCorrectingCodes.Codeword.hamming_distance}
\leanok
\statementsource{ada-coding-book}{pdf-7251cb8e4bc5-p018-b002}
\statementsource{philips}{pdf-24e94ad30862-p023-b007}
\statementsource{IEEE-TIT.23.157}{pdf-abed77da6b7a-p001-b004}
Hamming distance between two codewords, wrapping Mathlib's \texttt{hammingDist}.
\end{definition}

\begin{definition}[Distance predicate for a code]
\lean{ErrorCorrectingCodes.Codeword.distance}
\label{ErrorCorrectingCodes.Codeword.distance}
\leanok
\statementsource{ada-coding-book}{pdf-7251cb8e4bc5-p022-b008}
\statementsource{philips}{pdf-24e94ad30862-p025-b003}
\statementsource{IEEE-TIT.23.157}{pdf-abed77da6b7a-p001-b004}
\texttt{distance C d} holds when $d$ is the minimum Hamming distance of $C$: there
exist distinct codewords at distance exactly $d$, and no two distinct codewords
are closer than $d$.
\end{definition}

\begin{definition}[Rate]
\lean{ErrorCorrectingCodes.Codeword.rate}
\label{ErrorCorrectingCodes.Codeword.rate}
\leanok
\statementsource{ada-coding-book}{
  pdf-7251cb8e4bc5-p016-b003,
  pdf-7251cb8e4bc5-p017-b002
}
\statementsource{IEEE-TIT.23.157}{pdf-abed77da6b7a-p001-b004}
\[
R(C) = \frac{\log|C|}{n\log|\alpha|}.
\]
\end{definition}

\begin{definition}[Weight]
\lean{ErrorCorrectingCodes.Codeword.weight}
\label{ErrorCorrectingCodes.Codeword.weight}
\leanok
\statementsource{ada-coding-book}{pdf-7251cb8e4bc5-p027-b002}
\statementsource{philips}{pdf-24e94ad30862-p023-b009}
\statementsource{IEEE-TIT.23.157}{pdf-abed77da6b7a-p001-b004}
The Hamming weight of $c$ is $d(c, \mathbf{0})$.
\end{definition}

\begin{definition}[Maximal size at distance]
\lean{ErrorCorrectingCodes.Codeword.max_size}
\label{ErrorCorrectingCodes.Codeword.max_size}
\leanok
\statementsource{philips}{pdf-24e94ad30862-p025-b004}
\statementsource{IEEE-TIT.23.157}{pdf-abed77da6b7a-p001-b004}
\texttt{max\_size n d A} asserts existence of a code of size $A$ and minimum distance $d$
that is maximal among all such codes.
\end{definition}

\begin{definition}[Hamming ball]
\lean{ErrorCorrectingCodes.Codeword.hamming_ball}
\label{ErrorCorrectingCodes.Codeword.hamming_ball}
\leanok
\statementsource{ada-coding-book}{pdf-7251cb8e4bc5-p029-b006}
The Hamming ball of radius $l$ around $c$:
\[
B_l(c) = \{c' : d(c',c) \le l\},
\]
implemented as a \texttt{Finset}.
\end{definition}

%%%%%%%%%%%%%%%%%%%%%%%%%%%%%%%%%%%%%%%%%%%%%%%%%%%%%%%%%%%%%%%%%%%%%%%%%%%%%%%
\section{Basic lemma: distance vs.\ block length}

\begin{lemma}[Minimum distance is at most the block length]
\lean{ErrorCorrectingCodes.Codeword.dist_le_length}
\label{ErrorCorrectingCodes.Codeword.dist_le_length}
\leanok
\difficulty{3}
Let $C$ be a code of length $n$ over an alphabet $\alpha$ — that is, a finite set of
codewords, each a function $\mathrm{Fin}\,n \to \alpha$ — and let $d$ be a natural
number that is the minimum distance of $C$: there exist two distinct codewords of $C$ at
Hamming distance exactly $d$, and every pair of distinct codewords of $C$ is at Hamming
distance at least $d$. Then $d \le n$.
\end{lemma}

%%% added by the blueprint pipeline
\section{Additional declarations}

\begin{lemma}[Positivity of a natural number at least two]
\lean{ErrorCorrectingCodes.Codeword.natCast_pos_of_two_le}
\label{ErrorCorrectingCodes.Codeword.natCast_pos_of_two_le}
\leanok
\difficulty{1}
Let $q$ be a natural number with $q \ge 2$. Then its image in the real numbers satisfies
$0 < q$.
\end{lemma}

\begin{lemma}[One is less than a natural number at least two]
\lean{ErrorCorrectingCodes.Codeword.natCast_one_lt_of_two_le}
\label{ErrorCorrectingCodes.Codeword.natCast_one_lt_of_two_le}
\leanok
\difficulty{1}
For any natural number $q$ with $q \ge 2$, the image of $q$ under the canonical
embedding of the naturals into the reals satisfies $1 < q$.
\end{lemma}

\begin{lemma}[Positivity of $q-1$ for $q \ge 2$]
\lean{ErrorCorrectingCodes.Codeword.natCast_sub_one_pos_of_two_le}
\label{ErrorCorrectingCodes.Codeword.natCast_sub_one_pos_of_two_le}
\leanok
\uses{ErrorCorrectingCodes.Codeword.natCast_one_lt_of_two_le}
\difficulty{2}
Let $q$ be a natural number with $q \ge 2$. Then, regarded as a real number, $q - 1$ is
strictly positive: $0 < q - 1$.
\end{lemma}

\begin{lemma}[Integers at least two are not one as reals]
\lean{ErrorCorrectingCodes.Codeword.natCast_ne_one_of_two_le}
\label{ErrorCorrectingCodes.Codeword.natCast_ne_one_of_two_le}
\leanok
\uses{ErrorCorrectingCodes.Codeword.natCast_one_lt_of_two_le}
\difficulty{2}
Let $q$ be a natural number with $q \ge 2$. Then, regarded as a real number, $q \ne 1$.
\end{lemma}

\begin{lemma}[Positivity of one minus a subunital real]
\lean{ErrorCorrectingCodes.Codeword.one_sub_pos_of_lt_one}
\label{ErrorCorrectingCodes.Codeword.one_sub_pos_of_lt_one}
\leanok
\difficulty{1}
Let $p$ be a real number with $p < 1$. Then $1 - p > 0$.
\end{lemma}

\begin{lemma}[Positivity of $p(1-p)$ on the unit interval]
\lean{ErrorCorrectingCodes.Codeword.mul_one_sub_pos}
\label{ErrorCorrectingCodes.Codeword.mul_one_sub_pos}
\leanok
\uses{ErrorCorrectingCodes.Codeword.one_sub_pos_of_lt_one}
\difficulty{1}
Let $p$ be a real number with $0 < p < 1$. Then $p(1-p) > 0$.
\end{lemma}

\begin{lemma}[A rate bound below one]
\lean{ErrorCorrectingCodes.Codeword.lt_one_of_le_one_sub_inv}
\label{ErrorCorrectingCodes.Codeword.lt_one_of_le_one_sub_inv}
\leanok
\difficulty{2}
Let $q$ and $p$ be real numbers with $q > 0$. If $p \le 1 - 1/q$, then $p < 1$.
\end{lemma}

